% Abhakseite „Das kann ich“ – Zone nur im Gesamt (\mitzone)
%% TODO Ich-kann-Sätze: Platzhalter wie in den Aufgabendateien
\begin{abhakseite}
\ifdefined\mitzone
\abhakgruppe{Kennst du schon}
\abhak{1}{Termwerte berechnen}
\abhak{2}{Lineare Funktionen}
\abhak{3}{Maßstab und Einheiten}
\abhak{4}{Quadratische Gleichungen lösen}
\abhak{5}{Die Grundgraphen kennen}
\abhak{6}{Spiegeln im Koordinatensystem}
\abhak{7}{Termwerte berechnen – Fehler finden}
\abhak{8}{Termwerte berechnen – selbst rechnen}
\fi
\abhakgruppe{Einheit 1 · Funktionswert und Punkt}
\abhak{9}{Funktionswert und Punkt}
\abhak{10}{Nullstellen und Werte: Werte a mit vorgegebener Lösungsanzahl von f(x) = a über die lokalen Extremwerte am Graphen angeben}
\abhak{11}{Funktionswert und Punkt – Fehler finden, Begründen, Anwendung, Darstellungswechsel}
\abhakgruppe{Einheit 2 · Nullstellen}
\abhak{12}{Nullstellen}
\abhak{13}{Nullstellen und Werte: y-Achsenschnittpunkt angeben und Anzahl der Nullstellen aus Grenzverhalten und Tiefpunkt ohne Rechnung begründen}
\abhak{14}{Nullstellen – Fehler finden, Begründen, Anwendung}
\abhakgruppe{Einheit 3 · Einheit 3}
\abhak{15}{Definitionsbereich, Wertemenge und Schranken}
\abhak{16}{Definitionsbereich, Wertemenge und Schranken – Fehler finden, Begründen, Anwendung}
\abhakgruppe{Einheit 4 · Symmetrie}
\abhak{17}{Symmetrie}
\abhak{18}{Symmetrie – Fehler finden, Begründen}
\abhakgruppe{Einheit 5 · Transformationen}
\abhak{19}{Transformationen}
\abhak{20}{Transformationen – Fehler finden, Begründen, Darstellungswechsel, Anwendung}
\abhakgruppe{Einheit 6 · Graph und Term}
\abhak{21}{Graph und Term}
\abhak{22}{Graph und Term – Fehler finden, Begründen, Darstellungswechsel, Anwendung}
\end{abhakseite}
